% GENERATED FILE. Edit canonical lessons, then run npm run generate:books.
% canonical-source-hash: fnv1a64:c04f28034f5ddf4c
% canonical-lessons: ZH-C389-beiyong, ZH-C389-jieyue, ZH-C389-fengxian, ZH-C389-yifang, ZH-C389-daoqi

\chapter{In Case}
\label{ch:zh-in-case}
\hlchaptermodality{389}

\begin{chapteropening}
\textbf{By the end of this chapter:} \emph{I can talk about spare things, using less, risks, being ready for anything and payments falling due.}

Complete the last lesson of chapter 389: I can talk about spare things, using less, risks, being ready for anything and payments falling due.
\end{chapteropening}

\section[bèiyòng]{bèiyòng — spare (kept for when needed)}

\label{lesson:ZH-C389-beiyong}



\begin{quote}
Before the new one: say the Mandarin for to do ahead of time, then the Mandarin for a time zone.
\end{quote}

\subsection*{You'll want to know: bèiyòng}
\emph{bèiyòng} — \textquotedblleft{}spare (kept for when needed)\textquotedblright{}.

\begin{quote}
\emph{wànyī yàoshi diū le wǒmen hái yǒu bèiyòng de} — If the keys get lost, we've still got a spare.
\end{quote}

Its characters come later; for now it is a word for the ear.

\begin{cousinweb}
\emph{bèiyòng} is built from \emph{bèi} 'to provide against' and \zh{用} 'to use', both inherited from Middle Chinese, \emph{bijH} and \emph{yowngH} (Guangyun, in Baxter's transcription), which glosses the first 'complete' and 'to guard against'.
\end{cousinweb}

\begin{grammarlens}[title={What you've built}]
One more word for saying what would happen.
\end{grammarlens}

\subsection*{Your turn}
\begin{itemize}
\raggedright
  \item \emph{Say it:} \emph{bèiyòng}
  \item \emph{Say it:} \emph{bèiyòng}, once more
  \item \emph{Say it:} \emph{wànyī yàoshi diū le wǒmen hái yǒu bèiyòng de} — If the keys get lost, we've still got a spare.
  \item \emph{Recall:} say the Mandarin for to do ahead of time, then the Mandarin for a time zone, then say \emph{bèiyòng} again
\end{itemize}

\begin{tcolorbox}[breakable,colback=teal!4,colframe=teal!35!black,title={Before you move on}]
What does bèiyòng mean? (\textquotedblleft{}Spare (kept for when needed)\textquotedblright{}.) Say it once more.
\end{tcolorbox}
\section[\texorpdfstring{\zh{节约} (jiéyuē)}{节约 (jiéyuē)}]{\zh{节约} (jiéyuē) — to economise; to use sparingly}

\label{lesson:ZH-C389-jieyue}



\begin{quote}
Before the new one: say the Mandarin for a time zone, then the Mandarin for spare (kept for when needed).
\end{quote}

\subsection*{You'll want to know: \zh{节约}}
\textbf{\zh{节约}} — \emph{jiéyuē} — \textquotedblleft{}to economise; to use sparingly\textquotedblright{}.

\begin{quote}
\emph{rúguǒ dàjiā dōu jiéyuē yòng shuǐ jiù néng shěngqián} — If everyone uses water sparingly, we can all save money.
\end{quote}

\begin{cousinweb}
\textbf{\zh{节约}} (\emph{jiéyuē}) is built from \zh{节} 'to restrain' and \zh{约} 'sparing', both inherited from Middle Chinese, \emph{tset} and \emph{'jak} (Guangyun, in Baxter's transcription), which glosses the first 'to hold back' and the second 'thrifty'.
\end{cousinweb}

\begin{grammarlens}[title={What you've built}]
One more word for saying what would happen.
\end{grammarlens}

\subsection*{Your turn}
\begin{itemize}
\raggedright
  \item \emph{Say it:} \emph{jiéyuē}
  \item \emph{Say it:} \emph{jiéyuē}, once more
  \item \emph{Say it:} \emph{rúguǒ dàjiā dōu jiéyuē yòng shuǐ jiù néng shěngqián} — If everyone uses water sparingly, we can all save money.
  \item \emph{Recall:} say the Mandarin for a time zone, then the Mandarin for spare (kept for when needed), then say \emph{jiéyuē} again
\end{itemize}

\begin{tcolorbox}[breakable,colback=teal!4,colframe=teal!35!black,title={Before you move on}]
What does \zh{节约} mean? (\textquotedblleft{}To economise; to use sparingly\textquotedblright{}.) Say it once more.
\end{tcolorbox}
\section[fēngxiǎn]{fēngxiǎn — risk}

\label{lesson:ZH-C389-fengxian}



\begin{quote}
Before the new one: say the Mandarin for spare (kept for when needed), then the Mandarin for to economise.
\end{quote}

\subsection*{You'll want to know: fēngxiǎn}
\emph{fēngxiǎn} — \textquotedblleft{}risk\textquotedblright{}.

\begin{quote}
\emph{yàoshi fēngxiǎn tài dà wǒmen jiù bú zuò le} — If the risk is too great, we won't do it.
\end{quote}

Its characters come later; for now it is a word for the ear.

\begin{cousinweb}
\emph{fēngxiǎn} is built from \zh{风} 'the wind' and \emph{xiǎn} 'dangerous', both inherited from Middle Chinese, \emph{pjuwng} and \emph{xjemX} (Guangyun, in Baxter's transcription), which glosses the second 'perilous' and 'hard to pass'.
\end{cousinweb}

\begin{grammarlens}[title={What you've built}]
One more word for saying what would happen.
\end{grammarlens}

\subsection*{Your turn}
\begin{itemize}
\raggedright
  \item \emph{Say it:} \emph{fēngxiǎn}
  \item \emph{Say it:} \emph{fēngxiǎn}, once more
  \item \emph{Say it:} \emph{yàoshi fēngxiǎn tài dà wǒmen jiù bú zuò le} — If the risk is too great, we won't do it.
  \item \emph{Recall:} say the Mandarin for spare (kept for when needed), then the Mandarin for to economise, then say \emph{fēngxiǎn} again
\end{itemize}

\begin{tcolorbox}[breakable,colback=teal!4,colframe=teal!35!black,title={Before you move on}]
What does fēngxiǎn mean? (\textquotedblleft{}Risk\textquotedblright{}.) Say it once more.
\end{tcolorbox}
\section[yǐfáng]{yǐfáng — so as to guard against; as a precaution}

\label{lesson:ZH-C389-yifang}



\begin{quote}
Before the new one: say the Mandarin for to economise, then the Mandarin for a risk.
\end{quote}

\subsection*{You'll want to know: yǐfáng}
\emph{yǐfáng} — \textquotedblleft{}so as to guard against; as a precaution\textquotedblright{}.

\begin{quote}
\emph{dài yìdiǎnr qián yǐfáng wànyī} — Take a bit of money with you, in case.
\end{quote}

Its characters come later; for now it is a word for the ear.

\begin{cousinweb}
\emph{yǐfáng} is built from \zh{以} 'by means of' and \emph{fáng} 'to guard against', both inherited from Middle Chinese, \emph{yiX} and \emph{bjang} (Guangyun, in Baxter's transcription), which glosses the first 'to use' and the second 'to defend' and 'a dyke'.
\end{cousinweb}

\begin{grammarlens}[title={What you've built}]
One more word for saying what would happen.
\end{grammarlens}

\subsection*{Your turn}
\begin{itemize}
\raggedright
  \item \emph{Say it:} \emph{yǐfáng}
  \item \emph{Say it:} \emph{yǐfáng}, once more
  \item \emph{Say it:} \emph{dài yìdiǎnr qián yǐfáng wànyī} — Take a bit of money with you, in case.
  \item \emph{Recall:} say the Mandarin for to economise, then the Mandarin for a risk, then say \emph{yǐfáng} again
\end{itemize}

\begin{tcolorbox}[breakable,colback=teal!4,colframe=teal!35!black,title={Before you move on}]
What does yǐfáng mean? (\textquotedblleft{}So as to guard against; as a precaution\textquotedblright{}.) Say it once more.
\end{tcolorbox}
\section[\texorpdfstring{\zh{到期} (dàoqī)}{到期 (dàoqī)}]{\zh{到期} (dàoqī) — to fall due; to run out}

\label{lesson:ZH-C389-daoqi}



\begin{quote}
Before the new one: say the Mandarin for a risk, then the Mandarin for as a precaution.
\end{quote}

\subsection*{You'll want to know: \zh{到期}}
\textbf{\zh{到期}} — \emph{dàoqī} — \textquotedblleft{}to fall due; to run out\textquotedblright{}.

\begin{quote}
\emph{fángzū xià ge yuè jiù dàoqī le} — The rent falls due next month.
\end{quote}

\begin{cousinweb}
\textbf{\zh{到期}} (\emph{dàoqī}) is built from \zh{到} 'to arrive' and \zh{期} 'a set time', both inherited from Middle Chinese, \emph{tawH} and \emph{gi} (Guangyun, in Baxter's transcription), which glosses the second 'a limit' and 'an appointed meeting'. Middle Chinese \emph{gi} would regularly give a rising tone, qí; the level tone of qī is irregular.
\end{cousinweb}

\begin{grammarlens}[title={What you've built}]
One more word for saying what would happen.
\end{grammarlens}

\subsection*{Your turn}
\begin{itemize}
\raggedright
  \item \emph{Say it:} \emph{dàoqī}
  \item \emph{Say it:} \emph{dàoqī}, once more
  \item \emph{Say it:} \emph{fángzū xià ge yuè jiù dàoqī le} — The rent falls due next month.
  \item \emph{Recall:} say the Mandarin for a risk, then the Mandarin for as a precaution, then say \emph{dàoqī} again
\end{itemize}

\begin{tcolorbox}[breakable,colback=teal!4,colframe=teal!35!black,title={Before you move on}]
What does \zh{到期} mean? (\textquotedblleft{}To fall due; to run out\textquotedblright{}.) Say it once more.
\end{tcolorbox}
